\documentclass[notumble,10pt]{leaflet}

\usepackage[utf8]{inputenc}
\usepackage[spanish]{babel}
\usepackage[T1]{fontenc}
\usepackage{lmodern}
\usepackage{graphicx}
\usepackage{xcolor}
\usepackage{tcolorbox}
\usepackage{microtype}
% fontawesome5 no disponible - usar alternativas de texto
\usepackage{tikz}
\usetikzlibrary{babel, shapes.geometric, arrows.meta, positioning, shadows.blur, decorations.pathreplacing, calc, fadings}
\usepackage{pgfplots}
\pgfplotsset{compat=1.18}
\usepackage{amsmath}
\usepackage{enumitem}

% Definición de paleta de colores moderna y profesional
\definecolor{primary}{HTML}{1A365D}   % Azul oscuro corporativo
\definecolor{secondary}{HTML}{2B6CB0} % Azul brillante
\definecolor{accent}{HTML}{DD6B20}    % Naranja vibrante
\definecolor{bglight}{HTML}{F7FAFC}   % Gris muy claro (fondo)
\definecolor{textdark}{HTML}{2D3748}  % Gris oscuro (texto)

\pagecolor{bglight}
\color{textdark}

% Estilo de cajas para resaltar información
\tcbset{
    mybox/.style={
        colback=white,
        colframe=primary,
        boxrule=0.8pt,
        arc=4pt,
        left=5pt,right=5pt,top=5pt,bottom=5pt,
        fonttitle=\bfseries\large,
        coltitle=white,
        boxsep=2pt
    },
    accentbox/.style={
        colback=white,
        colframe=accent,
        boxrule=1pt,
        arc=4pt,
        left=5pt,right=5pt,top=5pt,bottom=5pt,
        fonttitle=\bfseries,
        coltitle=white,
    }
}

% Personalización de secciones
\usepackage{titlesec}
\titleformat{\section}{\Large\bfseries\color{primary}}{}{0em}{}[\vspace{-0.2em}\rule{\linewidth}{1pt}]
\titleformat{\subsection}{\large\bfseries\color{secondary}}{}{0em}{}

% --- TÍTULO (PORTADA - Panel 1) ---
\title{%
  \vspace{0.2cm}
  \normalsize\bfseries Universidad Nacional del Altiplano\\[0.05cm]
  \footnotesize Facultad de Ingeniería Mecánica, Eléctrica, Electrónica y Sistemas\\[0.05cm]
  \footnotesize Escuela Profesional de Ingeniería de Sistemas\\[0.3cm]
  \begin{tikzpicture}[baseline]
    \node[fill=primary, text=white, inner sep=10pt, rounded corners=8pt, align=center, text width=0.8\textwidth] {
      \Large\bfseries Algoritmos\\[0.3em] Meméticos
    };
  \end{tikzpicture}\\[0.3cm]
  \normalsize\color{secondary} Optimización Híbrida Inteligente\\[0.3cm]
  \small\color{primary} \textbf{Investigación de Operaciones}\\[0.3cm]
  \footnotesize Integrantes:\\[0.15cm]
  \scriptsize Barriga Nina Pablo Sebastian\\
  Condori Ticona Anthony\\
  Pari Valencia Jhon Daniel\\
  Jefry Edinson Ramos Castillo
}
\author{}
\date{}

\begin{document}
\maketitle
\thispagestyle{empty}

% Agregamos un gráfico en la portada
\begin{center}
\begin{tikzpicture}[scale=0.65, transform shape]
  % Fondo difuminado (Glow)
  \fill[secondary!15, blur shadow={shadow blur steps=10, shadow blur radius=1.5ex}] (0,0) circle (2.2cm);
  
  % Círculos orbitales concéntricos
  \draw[primary!30, dashed, thick] (0,0) circle (2.5cm);
  \draw[primary!50, thick] (0,0) circle (1.8cm);
  
  % Iconos centrales (ADN y Cerebro)
  \node[primary, scale=2.5] at (-0.6, 0.2) {\textbf{DNA}};
  \node[accent, scale=2.5] at (0.6, -0.2) {\textbf{Meme}};
  
  % Conexiones tipo red neuronal
  \draw[primary, thick, -{Stealth[length=2mm]}] (-0.3, 0.2) -- (0.2, -0.1);
  \draw[primary, thick, -{Stealth[length=2mm]}] (-0.6, -0.6) -- (-0.2, -0.2);
  \draw[primary, thick, -{Stealth[length=2mm]}] (0.6, 0.6) -- (0.2, 0.2);
  
  % Flechas circulares (Ciclo evolutivo)
  \foreach \angle in {0, 90, 180, 270} {
      \draw[->, >=Stealth, thick, accent, line width=1.2pt] 
          (\angle+10:2.5cm) arc (\angle+10:\angle+80:2.5cm);
  }
  
  % Nodos orbitales
  \node[circle, fill=primary, inner sep=3pt] at (45:2.5cm) {};
  \node[circle, fill=accent, inner sep=3pt] at (135:2.5cm) {};
  \node[circle, fill=primary, inner sep=3pt] at (225:2.5cm) {};
  \node[circle, fill=secondary, inner sep=3pt] at (315:2.5cm) {};
\end{tikzpicture}
\vspace{0.8cm}

\textit{Evolución Biológica + Evolución Cultural}
\end{center}

\newpage

% --- PANEL 2: INTRODUCCIÓN ---
\section{¿Qué son los Algoritmos Meméticos?}
Los \textbf{Algoritmos Meméticos (MA)} son métodos de optimización basados en poblaciones que combinan la evolución biológica (algoritmos genéticos) con el aprendizaje individual (búsqueda local).

El término "meme", acuñado por Richard Dawkins, representa una unidad de información cultural que se transmite, adapta y mejora durante la vida del individuo, a diferencia de los genes que son estáticos.

\begin{tcolorbox}[mybox, title=\textbf{!} El Concepto Clave]
Mientras que en un algoritmo genético los individuos nacen y mueren con las mismas características base, en un algoritmo memético los individuos \textit{aprenden} y \textit{mejoran} su adaptación antes de reproducirse.
\end{tcolorbox}

\subsection{Origen e Historia}
Fueron introducidos por Pablo Moscato en 1989. Surgieron como una alternativa para resolver problemas NP-duros donde los Algoritmos Genéticos tradicionales convergían prematuramente o eran ineficientes en la explotación local.

\begin{tcolorbox}[mybox, title=\textbf{*} Conocimiento del Problema]
Como establece el teorema de \textit{No-Free-Lunch} (Wolpert y Macready), el rendimiento de un algoritmo de búsqueda depende de la cantidad y calidad del conocimiento del problema que incorpora. Los MAs están diseñados intrínsecamente para explotar agresivamente dicho conocimiento a través de heurísticas.
\end{tcolorbox}

\newpage

% --- PANEL 3: ESTRUCTURA ---
\section{Arquitectura del Algoritmo}

La estructura básica de un MA es un ciclo evolutivo potenciado:

\begin{enumerate}[label=\textbf{\arabic*.}, leftmargin=*]
    \item \textbf{Inicialización:} Se genera una población inicial (usualmente con heurísticas).
    \item \textbf{Evaluación:} Se calcula el *fitness* de cada individuo.
    \item \textbf{Selección:} Se eligen los padres para la reproducción.
    \item \textbf{Recombinación (Crossover):} Se combinan genes para crear descendencia.
    \item \textbf{Mutación:} Se introducen pequeños cambios aleatorios.
    \item \textbf{\color{accent}Búsqueda Local:} ¡El paso memético! Los individuos exploran su vecindario para alcanzar un óptimo local.
    \item \textbf{Reemplazo:} La nueva generación sustituye a la anterior.
\end{enumerate}

\begin{tcolorbox}[accentbox, title=\textbf{>}\_ Pseudocódigo: Búsqueda Local, fontupper=\small\ttfamily]
\begin{flushleft}
Procedure Local-Search-Engine(current)\\
begin\\
\hspace*{0.5cm}repeat\\
\hspace*{1.0cm}new $\leftarrow$ GenerateNeighbor(current);\\
\hspace*{1.0cm}if (Fg(new) $\geq$ Fg(current)) then\\
\hspace*{1.5cm}current $\leftarrow$ new;\\
\hspace*{0.5cm}until TerminationCriterion();\\
\hspace*{0.5cm}return current;\\
end
\end{flushleft}
\end{tcolorbox}

\begin{tcolorbox}[accentbox, title=\textbf{<>} Ciclo de Población, fontupper=\small\ttfamily]
\begin{flushleft}
repeat\\
\hspace*{0.5cm}newpop $\leftarrow$ GenerateNewPopulation(pop);\\
\hspace*{0.5cm}pop $\leftarrow$ UpdatePopulation(pop, newpop);\\
\hspace*{0.5cm}if pop has converged then\\
\hspace*{1.0cm}pop $\leftarrow$ RestartPopulation(pop);\\
until TerminationCriterion();
\end{flushleft}
\end{tcolorbox}

\newpage

% --- PANEL 4: VENTAJAS Y APLICACIONES ---
\section{¿Por qué utilizarlos?}

\subsection{Ventajas Principales}
\begin{itemize}[leftmargin=*, itemsep=2pt]
    \item \textbf{+} \textbf{Sinergia:} Equilibrio perfecto entre exploración global (recombinación) y explotación local.
    \item \textbf{+} \textbf{Convergencia:} Alcanzan soluciones de alta calidad en menos generaciones que los AG puros.
    \item \textbf{+} \textbf{Flexibilidad:} Permiten integrar conocimiento específico del problema en la búsqueda local.
\end{itemize}

\subsection{Aplicaciones Reales}
Los MAs son el estado del arte para resolver problemas complejos de Ingeniería e Investigación de Operaciones:
\begin{itemize}[leftmargin=*]
    \item \textbf{Rutas de Vehículos (VRP):} Optimización logística y distribución.
    \item \textbf{Horarios (Timetabling):} Asignación de clases, turnos médicos y vuelos.
    \item \textbf{Machine Learning:} Entrenamiento de Redes Neuronales y selección de características.
    \item \textbf{Bioinformática:} Plegamiento de proteínas y alineamiento de secuencias.
    \item \textbf{Problemas Clásicos NP-Duros:} Particionamiento de Grafos (Graph Partitioning), Coloración de Grafos, Set Covering.
    \item \textbf{Problemas de Asignación:} Quadratic Assignment y Min Generalised Assignment.
\end{itemize}

\begin{tcolorbox}[mybox, title=\textbf{1} Rendimiento Comprobado]
Décadas de investigación han demostrado que los Algoritmos Meméticos superan consistentemente a sus contrapartes clásicas (AG puros o Simulated Annealing) en problemas fuertemente combinatorios.
\end{tcolorbox}

\newpage

% --- PANEL 5: DISEÑO Y DESAFÍOS ---
\section{Consideraciones de Diseño}

Para diseñar un Algoritmo Memético exitoso, se deben tomar decisiones críticas:

\begin{tcolorbox}[mybox, title=\textbf{=} Parámetros Clave]
\begin{itemize}[leftmargin=*]
    \item \textbf{Frecuencia:} ¿Con qué frecuencia se aplica la búsqueda local?
    \item \textbf{Intensidad:} ¿Cuántos pasos de búsqueda local se permiten por individuo?
    \item \textbf{Diversidad:} ¿Cómo evitar que los óptimos locales idénticos llenen la población?
\end{itemize}
\end{tcolorbox}

\subsection{Operadores de Recombinación}
A diferencia de los enfoques "ciegos" de los AG, los MAs suelen usar \textbf{Recombinación Heurística}. El conocimiento del problema se utiliza para guiar inteligentemente qué características (\textit{features}) de los padres se transmiten a la descendencia, asegurando viabilidad y alta calidad.

\subsection{El Equilibrio Computacional}
La búsqueda local consume recursos. El gran reto es el compromiso (\textit{trade-off}) entre el esfuerzo computacional invertido en la explotación local frente a la exploración evolutiva.

\begin{center}
\begin{tikzpicture}[scale=0.85, transform shape]
  % Fondo suave general
  \fill[secondary!5, rounded corners=3mm] (0,0) rectangle (4.5, 4.5);
  
  % Área de Balance Óptimo (Sweet Spot)
  \fill[inner color=accent!40, outer color=accent!0] (2.5, 2.5) circle (1.8cm);

  % Ejes
  \draw[thick, -{Stealth[length=2.5mm]}, textdark] (0,0) -- (5,0) node[anchor=north, font=\small\bfseries] {Exploración Global};
  \draw[thick, -{Stealth[length=2.5mm]}, textdark] (0,0) -- (0,5) node[anchor=east, font=\small\bfseries, rotate=90, yshift=4mm] {Explotación Local};
  
  % Etiquetas de Zonas
  \node[textdark!60, font=\scriptsize\bfseries, align=center] at (1.2, 4) {Convergencia\\Prematura};
  \node[textdark!60, font=\scriptsize\bfseries, align=center] at (4, 1) {Costo\\Excesivo};

  % Punto Óptimo Memético
  \fill[primary, blur shadow={shadow blur steps=5}] (2.5, 2.5) circle (4pt);
  \node[primary, font=\bfseries, fill=white, inner sep=2pt, rounded corners=2pt] at (2.5, 3) {Balance Óptimo MA};
  
  % Líneas guía elegantes
  \draw[densely dashed, thick, primary!50] (2.5, 0) -- (2.5, 2.5);
  \draw[densely dashed, thick, primary!50] (0, 2.5) -- (2.5, 2.5);
\end{tikzpicture}
\end{center}

\newpage

% --- PANEL 6: CONCLUSIÓN (CONTRAPORTADA) ---
\section{Conclusión}

Los Algoritmos Meméticos representan una de las aproximaciones más poderosas en la Inteligencia Computacional moderna. 

Al reconocer que la evolución biológica no es el único mecanismo de adaptación en la naturaleza, e incorporar la noción de evolución cultural (memes), logran un rendimiento superior en problemas de optimización combinatoria.

\vspace{1cm}
\begin{tcolorbox}[colback=primary!10, colframe=primary, arc=5pt]
\centering
\Large\textbf{-- Resumen}\\
\vspace{0.3cm}
\normalsize AG (Naturaleza) + BL (Aprendizaje) = MA (Excelencia)
\end{tcolorbox}

\vfill

\end{document}
